\documentclass[10pt,a4paper]{amsart}
\usepackage{amsmath,amssymb,amsthm,graphicx}
% Retained minimal preamble for the Stacks Project chapter flat.tex (More on Flatness).
% Verbatim excerpt of the chapter preamble: Xy-pic setup (native preamble lines 14-18) and the native
% theorem declarations with the mathematical macros (native preamble lines 154-217). Copied mechanically;
% nothing added, removed or redefined. The wrapper class and the amsmath/amssymb/amsthm/graphicx base
% packages are supplied by the assembler, so the class/externaldocument part of the native preamble is omitted.
% For commutative diagrams we use Xy-pic
\usepackage[all]{xy}

% We use 2cell for 2-commutative diagrams.
\xyoption{2cell}
% Theorem environments.
%
\theoremstyle{plain}
\newtheorem{theorem}[subsection]{Theorem}
\newtheorem{proposition}[subsection]{Proposition}
\newtheorem{lemma}[subsection]{Lemma}

\theoremstyle{definition}
\newtheorem{definition}[subsection]{Definition}
\newtheorem{example}[subsection]{Example}
\newtheorem{exercise}[subsection]{Exercise}
\newtheorem{situation}[subsection]{Situation}

\theoremstyle{remark}
\newtheorem{remark}[subsection]{Remark}
\newtheorem{remarks}[subsection]{Remarks}

\numberwithin{equation}{subsection}

% Macros
%
\def\lim{\mathop{\mathrm{lim}}\nolimits}
\def\colim{\mathop{\mathrm{colim}}\nolimits}
\def\Spec{\mathop{\mathrm{Spec}}}
\def\Hom{\mathop{\mathrm{Hom}}\nolimits}
\def\Ext{\mathop{\mathrm{Ext}}\nolimits}
\def\SheafHom{\mathop{\mathcal{H}\!\mathit{om}}\nolimits}
\def\SheafExt{\mathop{\mathcal{E}\!\mathit{xt}}\nolimits}
\def\Sch{\mathit{Sch}}
\def\Mor{\mathop{\mathrm{Mor}}\nolimits}
\def\Ob{\mathop{\mathrm{Ob}}\nolimits}
\def\Sh{\mathop{\mathit{Sh}}\nolimits}
\def\NL{\mathop{N\!L}\nolimits}
\def\CH{\mathop{\mathrm{CH}}\nolimits}
\def\proetale{{pro\text{-}\acute{e}tale}}
\def\etale{{\acute{e}tale}}
\def\QCoh{\mathit{QCoh}}
\def\Ker{\mathop{\mathrm{Ker}}}
\def\Im{\mathop{\mathrm{Im}}}
\def\Coker{\mathop{\mathrm{Coker}}}
\def\Coim{\mathop{\mathrm{Coim}}}

% Boxtimes
%
\DeclareMathSymbol{\boxtimes}{\mathbin}{AMSa}{"02}

%
% Macros for moduli stacks/spaces
%
\def\QCohstack{\mathcal{QC}\!\mathit{oh}}
\def\Cohstack{\mathcal{C}\!\mathit{oh}}
\def\Spacesstack{\mathcal{S}\!\mathit{paces}}
\def\Quotfunctor{\mathrm{Quot}}
\def\Hilbfunctor{\mathrm{Hilb}}
\def\Curvesstack{\mathcal{C}\!\mathit{urves}}
\def\Polarizedstack{\mathcal{P}\!\mathit{olarized}}
\def\Complexesstack{\mathcal{C}\!\mathit{omplexes}}
% \Pic is the operator that assigns to X its picard group, usage \Pic(X)
% \Picardstack_{X/B} denotes the Picard stack of X over B
% \Picardfunctor_{X/B} denotes the Picard functor of X over B
\def\Pic{\mathop{\mathrm{Pic}}\nolimits}
\def\Picardstack{\mathcal{P}\!\mathit{ic}}
\def\Picardfunctor{\mathrm{Pic}}
\def\Deformationcategory{\mathcal{D}\!\mathit{ef}}

% --- appended 2026-10-12 (Chief): the project own definitions of the reference/slogan/history
% display environments, copied mechanically from the shared Stacks preamble lines 8-12
% (\usepackage{verbatim} plus the three \newenvironment lines, which the chapter preamble keeps
% verbatim from the project preamble). Nothing added, removed or redefined: these environments
% carry no mathematical content and the project itself typesets them as comments. Without these
% three lines a counted statement whose own text uses them fails with Environment slogan undefined.
\usepackage{verbatim}
\newenvironment{reference}{\comment}{\endcomment}
\newenvironment{slogan}{\comment}{\endcomment}
\newenvironment{history}{\comment}{\endcomment}
\usepackage{cleveref}
\crefname{originaltheorem}{original theorem}{original theorems}
\Crefname{originaltheorem}{Original theorem}{Original theorems}
\makeatletter
\newlabel{theorem-flattening-map}{{05PF}{source}{}{sourcepointer-1}{}}
\newlabel{theorem-flattening-map@cref}{{[originaltheorem][0][]05PF}{[1][0][]source}{}{}{}}
\makeatother
\title{Weil restriction along a flat pure morphism of finite presentation is representable by a closed immersion, of finite presentation when the subscheme is}
\author{The Stacks Project Authors}
\date{stacks-project:flat.tex:local-snapshot-20261011; GNU FDL 1.2 or later}
\begin{document}
\maketitle
\noindent\textit{Editorial note.} All delivered mathematical text is the author's own, copied from the pinned Stacks Project chapter file flat.tex; no statement or proof step was rewritten, reordered or removed.\par\medskip
\noindent\textit{Editorial note.} Preamble provenance: the chapter's own retained minimal preamble (sources/stacks-flat/retained-preamble.tex, lines 1-85) is a verbatim excerpt of the native Xy-pic setup and of the native theorem declarations and mathematical macros, followed by the project's own reference/slogan/history display-environment definitions copied mechanically from the shared preamble; the wrapper class and the amsmath/amssymb/amsthm/graphicx base packages are supplied by the assembler, and extra\_packages is left empty because everything needed is already in the retained preamble. No macro was invented or redefined.\par\medskip
\noindent\textit{Editorial note.} Same-chapter pointers are located at their \textbackslash{}label lines inside the pinned chapter file; cross-chapter pointers are located in the project's own tag list staged in the snapshot (sources/stacks-chow/tags.txt), at the verbatim lines carrying their canonical tags.\par\medskip
\noindent\textit{Allowed standard prerequisites (editorial).}\begin{itemize}
\item The flattening map theorem, representing the locus where a map of modules becomes an isomorphism. Same chapter flat.tex, label theorem-flattening-map (invoked at proof line 6777). Applied to the surjection \$u : \textbackslash{}mathcal\{O\}\_X \textbackslash{}to \textbackslash{}mathcal\{O\}\_X/\textbackslash{}mathcal\{I\}\$, it represents by a closed immersion \$Z \textbackslash{}to S\$ the \$T\$ for which \$u\_T\$ is an isomorphism, which is exactly the Weil restriction functor \$F\$.
\item The ideal sheaf of a closed subscheme and the associated surjection. Same chapter flat.tex, the ideal sheaf \$\textbackslash{}mathcal\{I\}\$ of \$Y\$ in \$X\$ (used at proof lines 6772-6773). It converts the condition that \$Y\_T \textbackslash{}to X\_T\$ be an isomorphism into the condition that the pullback \$u\_T\$ of the surjection be an isomorphism.
\end{itemize}\medskip
\ifdefined\hypertarget\hypertarget{sourcepointer-1}{}\fi
\noindent\textit{Source reference.} Original theorem 05PF: sources/stacks-flat/pinned-original-source.tex, source lines 6671--6671. Navigation pointer for \textbackslash{}ref\{theorem-flattening-map\} used inside the counted span (line 6777); the argument that consumes it is retained in the span.\par\medskip
\begin{lemma}
\label{lemma-Weil-restriction-closed-subschemes}
Let $f:X\to S$ be a morphism of schemes which is of finite presentation,
flat, and pure. Let $Y$ be a closed subscheme of $X$. Let $F=f_*Y$ be the
Weil restriction functor of $Y$ along $f$, defined by
$$
F : (\Sch/S)^{opp} \to \textit{Sets}, \quad
T \mapsto
\left\{
\begin{matrix}
\{*\} & \text{if} & Y_T\to X_T \text{ is an isomorphism, }\\
\emptyset & \text{else.} &
\end{matrix}
\right.
$$
Then $F$ is representable by a closed immersion $Z\to S$. Moreover
$Z\to S$ is of finite presentation if $Y\to S$ is.
\end{lemma}

\begin{proof}
Let $\mathcal{I}$ be the ideal sheaf defining $Y$ in $X$ and let
$u:\mathcal{O}_X\to\mathcal{O}_X/\mathcal{I}$ be the surjection.
Then for an $S$-scheme $T$, the closed immersion $Y_T\to X_T$
is an isomorphism if and only if $u_T$ is an isomorphism. Hence
the result follows from
Theorem \ref{theorem-flattening-map}.
\end{proof}

\end{document}
